% Kuis 2 --- Logika Predikat
% Satu soal PG per bahasan utama di ppt/Logika_Predikat.tex.
% Skenario berbeda dari tugas/tugas2.tex dan UAS/uas.tex.
\documentclass[12pt,a4paper]{article}

\usepackage[utf8]{inputenc}
\usepackage[T1]{fontenc}
\usepackage[indonesian]{babel}

\usepackage{geometry}
\geometry{margin=2cm}

\usepackage{lmodern}
\usepackage{amsmath,amssymb}
\usepackage{enumitem}
\usepackage{tabularx}
\usepackage{microtype}

\usepackage{fancyhdr}
\usepackage{lastpage}

\setlength{\headheight}{14pt}
\addtolength{\topmargin}{-2pt}
\usepackage[hidelinks,breaklinks]{hyperref}

\newcommand{\JudulKuis}{Kuis 2}
\newcommand{\KodeMK}{TIF-xxx}
\newcommand{\NamaMK}{Logika Matematika}
\newcommand{\MateriKuis}{Logika Predikat}
\newcommand{\WaktuKuis}{90 menit}
\newcommand{\NilaiMaksimal}{100}
\newcommand{\Prodi}{Teknik Informatika}
\newcommand{\Fakultas}{Fakultas Teknologi Informasi}
\newcommand{\Institusi}{Universitas Bale Bandung}
\newcommand{\DosenPengampu}{[Nama Dosen Pengampu]}
\newcommand{\TanggalKuis}{[DD Bulan YYYY]}

\pagestyle{fancy}
\fancyhf{}
\fancyhead[L]{\small \JudulKuis\ --- \NamaMK}
\fancyhead[R]{\small \MateriKuis}
\fancyfoot[C]{\small Halaman \thepage\ dari \pageref*{LastPage}}
\renewcommand{\headrulewidth}{0.4pt}
\renewcommand{\footrulewidth}{0pt}

\setlength{\parindent}{0pt}
\setlength{\parskip}{0.5em}

\newlist{pilihan}{enumerate}{1}
\setlist[pilihan]{label=\textbf{\Alph*.}, leftmargin=2.2em, align=left, itemsep=0.35em, topsep=0.35em}

\newcommand{\kuncijawaban}[2]{%
  \par\vspace{0.25em}%
  {\small
  \textit{\textbf{Kunci jawaban:} #1}\par
  \textit{\textbf{Penjelasan:} #2}%
  \par\vspace{0.55em}}%
}

\newcounter{soalkuis}
\newcommand{\soalkuis}{%
  \refstepcounter{soalkuis}%
  \par\vspace{0.6ex plus 0.15ex}%
  \noindent\textbf{Soal \thesoalkuis.}\quad
}

\begin{document}

\begin{center}
  {\Large\bfseries \JudulKuis}\\[0.3em]
  {\large \NamaMK\ (\KodeMK)}\\[0.25em]
  {\normalsize Materi: \MateriKuis}\\[0.35em]
  {\normalsize\textit{(Lembar kunci jawaban dan pembahasan)}}\\[0.8em]
\end{center}

\noindent
\begin{tabularx}{\linewidth}{@{}p{4.5cm} @{\hspace{0.35em}:\hspace{0.65em}} >{\raggedright\arraybackslash}X@{}}
  Nama Mahasiswa        & \rule{0.72\linewidth}{0.4pt} \\
  Nomor Induk Mahasiswa & \rule{0.72\linewidth}{0.4pt} \\
  Kelas / Kelompok      & \rule{0.72\linewidth}{0.4pt} \\
  Program Studi         & \Prodi \\
  Fakultas              & \Fakultas \\
  Institusi             & \Institusi \\
  Dosen Pengampu        & \DosenPengampu \\
  Tanggal kuis          & \TanggalKuis \\
  Waktu pengerjaan      & \WaktuKuis \\
  Nilai maksimal        & \NilaiMaksimal \\
\end{tabularx}

\vspace{0.8em}
\hrule
\vspace{0.8em}

\section*{Petunjuk}
\begin{itemize}[leftmargin=*, itemsep=0.25em]
  \item Kuis ini terdiri atas \textbf{18 soal pilihan ganda} (satu jawaban benar per soal), mencakup bahasan utama \textit{Logika Predikat} pada materi perkuliahan: motivasi kalkulus predikat, simbol, term, proposisi, kalimat, ekspresi, subterm/subkalimat, representasi bahasa alami, variabel bebas/terikat, kalimat tertutup, simbol bebas, interpretasi, arti kalimat, aturan semantik kuantor, interpretasi diperluas, kecocokan, dan validitas.
  \item Isi \textbf{nama, NIM, dan kelas} pada bagian identitas di atas sebelum mengerjakan.
  \item Bacalah teks soal dan label bahasan (cetak miring di awal soal) untuk memastikan konteks pertanyaan.
  \item Pilih \textbf{satu} jawaban terbaik; lingkari atau tulis huruf pilihan (A, B, C, atau D) dengan jelas pada lembar jawaban.
  \item Gunakan pengertian predikat, kuantor \(\forall\) dan \(\exists\), variabel bebas/terikat, serta aturan semantik sesuai slide kuliah.
  \item Perangkat elektronik, catatan, dan komunikasi antar peserta \textbf{tidak diperbolehkan} kecuali diizinkan pengawas.
  \item Lembar kuis \textbf{wajib dikumpulkan} kepada pengawas segera setelah waktu \WaktuKuis\ habis; keterlambatan mengikuti aturan kelas.
\end{itemize}

\section*{Soal}

% --- 1 Pendahuluan ---
\soalkuis
\textit{(Pendahuluan / Motivasi)}\\
Kalimat ``Semua burung dapat terbang'' sulit dinyatakan secara memadai hanya dengan kalkulus proposisi karena \ldots
\begin{pilihan}
  \item Kalimat tersebut bukan pernyataan sama sekali
  \item Kalimat tersebut mengandung kuantisasi atas objek beserta sifat/relasinya
  \item Kalkulus proposisi tidak mengenal simbol \(\neg\)
  \item Nilai kebenaran selalu indeterminate
\end{pilihan}
\kuncijawaban{B}{Kalkulus proposisi tidak merepresentasikan objek dan kuantor (``semua''/``ada''); untuk itu diperlukan kalkulus predikat.}

% --- 2 Definisi Simbol ---
\soalkuis
\textit{(Definisi Simbol)}\\
Manakah yang termasuk \textbf{simbol predikat} dalam kalkulus predikat?
\begin{pilihan}
  \item \(x, y, z\)
  \item \(a, b, c\)
  \item \(P, Q, R\) (dengan arity tertentu) yang menyatakan sifat atau relasi
  \item \(\mathit{true}\) dan \(\mathit{false}\) saja
\end{pilihan}
\kuncijawaban{C}{Simbol predikat menyatakan sifat atau relasi antar objek dan memiliki arity; variabel adalah \(x,y,\ldots\), konstanta \(a,b,\ldots\).}

% --- 3 Definisi Term ---
\soalkuis
\textit{(Definisi Term)}\\
Asumsikan \(a\) konstanta, \(x\) variabel, dan \(f\) fungsi unary.
Manakah yang merupakan \textbf{term}?
\begin{pilihan}
  \item \(P(a)\)
  \item \(f(a)\)
  \item \(\forall x\,P(x)\)
  \item \(P(x) \wedge Q(a)\)
\end{pilihan}
\kuncijawaban{B}{Term menyatakan objek: konstanta, variabel, atau aplikasi fungsi atas term. \(P(a)\) adalah proposisi/kalimat atomik, bukan term.}

% --- 4 Definisi Proposisi ---
\soalkuis
\textit{(Definisi Proposisi)}\\
Jika \(t_1,\ldots,t_n\) adalah term dan \(P\) predikat \(n\)-ary, maka \(P(t_1,\ldots,t_n)\) adalah \ldots
\begin{pilihan}
  \item Fungsi
  \item Konstanta
  \item Proposisi (atomik) yang merepresentasikan relasi antar objek
  \item Kuantor
\end{pilihan}
\kuncijawaban{C}{Menurut definisi, predikat yang diterapkan pada term menghasilkan proposisi yang menyatakan relasi antar objek.}

% --- 5 Definisi Kalimat ---
\soalkuis
\textit{(Definisi Kalimat)}\\
Jika \(A\) adalah kalimat dan \(x\) variabel, maka \((\forall x)\,A\) adalah \ldots
\begin{pilihan}
  \item Term
  \item Fungsi
  \item Kalimat
  \item Bukan ekspresi yang terbentuk dengan baik
\end{pilihan}
\kuncijawaban{C}{Aturan pembentukan kalimat mencakup kuantifikasi: jika \(A\) kalimat maka \((\forall x)\,A\) dan \((\exists x)\,A\) juga kalimat.}

% --- 6 Definisi Ekspresi ---
\soalkuis
\textit{(Definisi Ekspresi)}\\
Dalam kalkulus predikat, suatu \textbf{ekspresi} dapat berupa \ldots
\begin{pilihan}
  \item Hanya kuantor
  \item Hanya konstanta kebenaran
  \item Kalimat atau term
  \item Hanya tabel kebenaran
\end{pilihan}
\kuncijawaban{C}{Materi mendefinisikan ekspresi sebagai kalimat atau term (misalnya \(x\), \(f(x,y)\), atau \((\exists x)\,P(x)\)).}

% --- 7 Subterm / Subkalimat ---
\soalkuis
\textit{(Subterm / Subkalimat)}\\
Pada ekspresi \(E :\; R\bigl(g(b), h(w)\bigr)\), manakah himpunan yang memuat \textbf{semua subterm} yang jelas tampak?
\begin{pilihan}
  \item Hanya \(R(g(b), h(w))\)
  \item \(b\), \(w\), \(g(b)\), \(h(w)\), dan \(R(g(b), h(w))\) sebagai term/subterm terkait objek (bukan predikat sebagai term)
  \item \(b\), \(w\), \(g(b)\), \(h(w)\)
  \item Hanya \(R\)
\end{pilihan}
\kuncijawaban{C}{Subterm adalah term antara yang membangun ekspresi: konstanta/variabel \(b,w\) dan aplikasi fungsi \(g(b)\), \(h(w)\). Predikat \(R(\ldots)\) adalah proposisi/kalimat, bukan term.}

% --- 8 Representasi: semua ---
\soalkuis
\textit{(Representasi Kalimat --- ``Semua'')}\\
Misalkan \(B(x)\): ``\(x\) adalah kelelawar'', \(F(x)\): ``\(x\) adalah lalat'', \(E(y,x)\): ``\(y\) memakan \(x\)''.
Terjemahan terbaik untuk ``Hanya kelelawar yang memakan lalat'' adalah \ldots
\begin{pilihan}
  \item \(\forall x\,(F(x) \rightarrow \forall y\,(E(y,x) \rightarrow B(y)))\)
  \item \(\forall x\,(B(x) \rightarrow F(x))\)
  \item \(\exists x\,(B(x) \wedge F(x))\)
  \item \(\forall x\,(F(x) \wedge B(x))\)
\end{pilihan}
\kuncijawaban{A}{``Hanya kelelawar memakan lalat'' berarti: untuk setiap lalat \(x\), setiap pemakan \(y\) dari \(x\) adalah kelelawar; ini cocok dengan \(\forall x\,(F(x)\rightarrow\forall y\,(E(y,x)\rightarrow B(y)))\).}

% --- 9 Representasi: ada / tidak semua ---
\soalkuis
\textit{(Representasi Kalimat --- ``Tidak semua'')}\\
Pernyataan ``Tidak semua bilangan bulat adalah genap'' paling tepat dinyatakan sebagai \ldots
(gunakan \(Z(x)\): ``\(x\) bilangan bulat'', \(G(x)\): ``\(x\) genap'')
\begin{pilihan}
  \item \(\forall x\,(Z(x) \rightarrow G(x))\)
  \item \(\exists x\,(Z(x) \wedge \neg G(x))\)
  \item \(\neg\exists x\,(Z(x) \wedge G(x))\)
  \item \(\forall x\,(Z(x) \wedge G(x))\)
\end{pilihan}
\kuncijawaban{B}{``Tidak semua bilangan bulat genap'' ekuivalen dengan ``ada bilangan bulat yang tidak genap'': \(\exists x\,(Z(x)\wedge\neg G(x))\).}

% --- 10 Variabel bebas/terikat ---
\soalkuis
\textit{(Variabel Bebas atau Terikat)}\\
Pada formula \(\forall x\,P(x,y)\), variabel yang \textbf{bebas} adalah \ldots
\begin{pilihan}
  \item Hanya \(x\)
  \item Hanya \(y\)
  \item \(x\) dan \(y\)
  \item Tidak ada; keduanya terikat
\end{pilihan}
\kuncijawaban{B}{Kuantor \(\forall x\) mengikat kemunculan \(x\); \(y\) tidak berada dalam cakupan kuantor atas \(y\), sehingga \(y\) bebas.}

% --- 11 Kalimat tertutup ---
\soalkuis
\textit{(Kalimat Tertutup)}\\
Manakah yang merupakan \textbf{kalimat tertutup}?
\begin{pilihan}
  \item \(\forall x\,P(x,y)\)
  \item \(P(a,b) \wedge Q(x)\)
  \item \(\forall x\,\exists y\,R(x,y)\)
  \item \(\exists y\,R(x,y)\)
\end{pilihan}
\kuncijawaban{C}{Kalimat tertutup tidak memiliki variabel bebas. Pada \(\forall x\,\exists y\,R(x,y)\) semua kemunculan \(x\) dan \(y\) terikat.}

% --- 12 Simbol bebas ---
\soalkuis
\textit{(Simbol Bebas)}\\
Simbol bebas dari suatu ekspresi \(A\) mencakup \ldots
\begin{pilihan}
  \item Hanya kuantor yang muncul di \(A\)
  \item Variabel bebas, konstanta, simbol fungsi, dan simbol predikat dari \(A\)
  \item Hanya variabel terikat
  \item Hanya domain interpretasi
\end{pilihan}
\kuncijawaban{B}{Menurut materi, simbol bebas meliputi variabel bebas serta semua konstanta, fungsi, dan predikat yang muncul pada ekspresi.}

% --- 13 Interpretasi ---
\soalkuis
\textit{(Interpretasi)}\\
Sebuah interpretasi \(I\) untuk kalimat predikat harus menyediakan \ldots
\begin{pilihan}
  \item Hanya tabel kebenaran proposisional
  \item Domain tak kosong serta pemetaan nilai untuk setiap simbol bebas yang relevan
  \item Hanya daftar kuantor
  \item Hanya nama mahasiswa
\end{pilihan}
\kuncijawaban{B}{Interpretasi mendefinisikan domain objek tak kosong dan memberi nilai pada konstanta, variabel bebas, fungsi, serta predikat (simbol bebas).}

% --- 14 Arti kalimat ---
\soalkuis
\textit{(Arti Kalimat)}\\
Misalkan domain \(D=\{0,1,2\}\), \(a_I=0\), dan \(P_I=\{0,2\}\) (ekstensi predikat unary \(P\)).
Nilai dari \((\exists x)\,P(x)\) berdasarkan \(I\) adalah \ldots
\begin{pilihan}
  \item Salah, karena \(1\notin P_I\)
  \item Benar, karena ada elemen (misalnya \(0\)) di \(D\) yang memenuhi \(P\)
  \item Tidak dapat ditentukan tanpa kuantor \(\forall\)
  \item Selalu salah pada domain berhingga
\end{pilihan}
\kuncijawaban{B}{Aturan FOR SOME: kalimat benar jika ada setidaknya satu elemen domain yang membuat \(P(x)\) benar; di sini \(0\) (dan \(2\)) memenuhi.}

% --- 15 Aturan semantik kuantor ---
\soalkuis
\textit{(Aturan Semantik untuk Kuantor)}\\
Kalimat \((\forall x)\,A\) bernilai \textbf{salah} berdasarkan interpretasi \(I\) dengan domain \(D\) jika dan hanya jika \ldots
\begin{pilihan}
  \item Untuk setiap \(d\in D\), \(A\) benar pada \(\langle x\leftarrow d\rangle\circ I\)
  \item Ada \(d\in D\) sedemikian sehingga \(A\) salah pada \(\langle x\leftarrow d\rangle\circ I\)
  \item \(A\) tidak mengandung variabel \(x\)
  \item Domain \(D\) kosong
\end{pilihan}
\kuncijawaban{B}{Menurut aturan FOR ALL: \(\forall x\,A\) salah tepat ketika ada elemen \(d\) yang membuat \(A\) salah di bawah interpretasi diperluas \(\langle x\leftarrow d\rangle\circ I\).}

% --- 16 Interpretasi diperluas ---
\soalkuis
\textit{(Interpretasi yang Diperluas)}\\
Notasi \(\langle x \leftarrow d\rangle \circ I\) berarti \ldots
\begin{pilihan}
  \item Menghapus seluruh domain \(I\)
  \item Interpretasi yang sama dengan \(I\), kecuali variabel \(x\) diberi nilai elemen \(d\)
  \item Mengganti semua predikat menjadi konstanta
  \item Menegasikan seluruh kalimat
\end{pilihan}
\kuncijawaban{B}{Interpretasi diperluas mempertahankan domain dan pemetaan lain dari \(I\), tetapi menetapkan (atau menimpa) nilai variabel \(x\) menjadi \(d\).}

% --- 17 Kecocokan ---
\soalkuis
\textit{(Kecocokan)}\\
Misalkan \(I\) dan \(J\) memberi nilai yang sama untuk konstanta \(a\) dan fungsi \(f\), tetapi \(I\) memberi \(x\mapsto 1\) sedangkan \(J\) memberi \(x\mapsto 3\).
Kesimpulan yang tepat adalah \ldots
\begin{pilihan}
  \item \(I\) dan \(J\) cocok untuk ekspresi \(f(a)\)
  \item \(I\) dan \(J\) cocok untuk ekspresi \(f(x)\)
  \item \(I\) dan \(J\) selalu cocok untuk segala ekspresi
  \item \(I\) dan \(J\) tidak cocok untuk konstanta \(a\)
\end{pilihan}
\kuncijawaban{A}{Kecocokan mensyaratkan nilai sama pada simbol yang relevan. Untuk \(f(a)\) hanya \(f\) dan \(a\) dibutuhkan (keduanya sama); untuk \(f(x)\) nilai \(x\) berbeda sehingga tidak cocok.}

% --- 18 Validitas ---
\soalkuis
\textit{(Validitas)}\\
Manakah pernyataan yang benar tentang validitas dalam kalkulus predikat?
\begin{pilihan}
  \item Setiap kalimat dengan \(\exists\) otomatis valid
  \item Kalimat tertutup \(A\) valid jika \(A\) bernilai benar pada \textbf{setiap} interpretasi untuk \(A\)
  \item Validitas hanya ditentukan dari satu baris tabel kebenaran proposisi
  \item Kalimat dengan variabel bebas selalu valid
\end{pilihan}
\kuncijawaban{B}{Validitas didefinisikan untuk kalimat tertutup: \(A\) valid bila benar berdasarkan setiap interpretasi. Untuk menunjukkan tidak valid cukup satu interpretasi kontra.}

\end{document}
